%% ============================================================
%%  Chapter 20 --- Parallel Execution and Fan-Out/Fan-In
%%  Production-Grade Agentic AI: LangGraph + Python Frameworks
%%  Dependencies: langgraph >= 0.2.0,
%%                langchain-litellm >= 0.3.0,
%%                pydantic >= 2.7.0,
%%                langchain-core >= 0.3.0
%% ============================================================

\chapterquote{Serialisation is not safety. It is the refusal to admit that the work was independent all along.}{anon}

The production failure this chapter prevents is serial latency
masquerading as a design choice. The supervisor pattern of Chapter~19
activates one worker per routing cycle: the LLM decides which worker
runs, that worker completes, control returns to the supervisor, and the
cycle repeats. For a research task that requires five independent
information sources, this means five sequential LLM-plus-retrieval round
trips. At three seconds per round trip the minimum wall-clock cost is
fifteen seconds, regardless of how fast each individual call is. The
failure is not occasional---it is the guaranteed lower bound of every
invocation. Deferring parallelism until ``performance becomes a problem''
misunderstands the cost: the graph is structurally incapable of doing
better because its topology serialises work that has no data dependency.

The solution has five layers. The first is a \emph{reducer layer}: every
field in graph state that will receive writes from parallel workers must
declare an accumulation function through \code{Annotated}; the
accumulation semantics determine correctness, and they must be chosen
before any parallel code is written. The second is a \emph{fan-out layer}:
a \code{decompose\_node} calls an LLM to split the query into at most
five sub-queries and writes them to state; a pure routing function
\code{fan\_out\_router} reads those sub-queries and returns a
\code{list[Send]}, one per sub-query, which LangGraph dispatches as
concurrent tasks. The third is a \emph{worker layer}: each dispatched
\code{search\_worker} task runs independently, calling two retrieval
backends concurrently with \code{asyncio.gather}; a \code{RetryPolicy}
registered at node-definition time handles transient infrastructure
failures at the LangGraph level without any code inside the worker. The
fourth is a \emph{fan-in and aggregation layer}: LangGraph's Pregel
scheduler waits for all parallel tasks to complete, then advances to
\code{aggregate\_node}, which receives the fully merged
\code{worker\_results} field and applies Reciprocal Rank Fusion across
all workers' results. The fifth is a \emph{synthesis layer}: a separate
LLM node distinct from the deterministic aggregation step produces the
final answer, keeping the two functions independently cacheable.

The sections are ordered by dependency.
Section~\ref{sec:reducers-parallel} establishes reducer semantics first
because every subsequent section produces or consumes a field whose
correctness depends on the reducer choice.
Section~\ref{sec:send-fanout} introduces \code{Send} and the
decompose-then-route pattern next because the worker's input type cannot
be specified before the fan-out contract is understood.
Section~\ref{sec:partial-failure} adds the worker implementation and its
two-level error containment after the worker's state type is established.
Section~\ref{sec:map-reduce-topo} shows the complete map-reduce topology
after all four node types exist.
Section~\ref{sec:full-parallel} assembles the production graph without
introducing new concepts.

This chapter introduces seven terms defined in Table~\ref{tab:ch20-terms}.

\begin{table}[h]
\centering
\caption{Domain terms introduced in Chapter~20}
\label{tab:ch20-terms}
\begin{tabularx}{\textwidth}{lX}
\toprule
\textbf{Term} & \textbf{Definition} \\
\midrule
\code{Send} &
  \code{from langgraph.types import Send}; a two-field object
  \code{Send(node,~arg)} where \code{node} is the name of the target
  node and \code{arg} is the entire state dict that node will receive;
  returned as a \code{list} from a conditional edge routing function
  to spawn multiple concurrent node instances, each with its own
  independent per-instance state. \\
Fan-out router &
  A pure synchronous function registered as the second argument to
  \code{add\_conditional\_edges}; reads parent graph state and returns
  \code{list[Send]}; never registered with \code{add\_node} and never
  performs I/O. \\
\code{WorkerInput} &
  A \code{TypedDict} with fields \code{sub\_query:~str} and
  \code{worker\_id:~int}; the type annotation for
  \code{search\_worker}'s \code{state} argument; matches the plain
  \code{dict} literal passed as \code{Send}'s \code{arg}. \\
\code{operator.add} reducer &
  \code{Annotated[list[T],~operator.add]} on a state field instructs
  LangGraph to accumulate parallel worker outputs by list concatenation;
  every completed worker's return list is appended to the accumulated
  list, preserving all entries without deduplication. \\
\code{RetryPolicy} &
  \code{from langgraph.types import RetryPolicy}; a \code{NamedTuple}
  passed as the \code{retry} keyword argument to \code{add\_node};
  key fields: \code{max\_attempts} (default~3, counting the initial
  attempt), \code{initial\_interval} (0.5~s), \code{backoff\_factor}
  (2.0), \code{max\_interval} (128~s), \code{jitter} (True),
  \code{retry\_on} (exception type, sequence, or callable). \\
\code{CachePolicy} &
  \code{from langgraph.types import CachePolicy}; a dataclass passed
  as the \code{cache} keyword argument to \code{add\_node}; fields:
  \code{key\_func} (defaults to pickle-hash of node input) and
  \code{ttl} (seconds; \code{None} = never expires within the
  process). \\
Pregel fan-in &
  LangGraph's Pregel-based scheduler accumulates reducer updates from
  each completed \code{Send} task and advances the graph to the next
  node only after \emph{all} tasks in the current batch have either
  completed or been permanently failed following \code{RetryPolicy}
  exhaustion. \\
\bottomrule
\end{tabularx}
\end{table}

\medskip
\noindent\textbf{Source layout for this chapter:}
\begin{verbatim}
chapter20/
  types.py        # §20.1 — SearchResult (all 5 fields incl. worker_id)
  state.py        # §20.1 — WorkerInput, ResearchState,
                  #          operator.add reducer, sub_queries field
  fan_out.py      # §20.2 — decompose_node (async, LLM),
                  #          fan_out_router (sync, pure routing),
                  #          MAX_WORKERS, SubQueryList
  worker.py       # §20.3 — search_worker with asyncio.gather +
                  #          return_exceptions=True
  backends.py     # §20.3 — qdrant_search, bm25_search stubs
  aggregate.py    # §20.4 — aggregate_node (RRF + dedup),
                  #          synthesise_node
  graph.py        # §20.5 — build_research_graph(), run_research()
  test_graph.py   # §20.5 — integration test with backend mocks
\end{verbatim}


%% ------------------------------------------------------------
\section{The Reducer Determines What Parallel Means for Each Field}
\label{sec:reducers-parallel}
%% ------------------------------------------------------------

The design problem is correctness under concurrent writes to shared
state. When LangGraph dispatches five \code{search\_worker} tasks in
parallel, all five eventually return a partial state update whose key is
\code{worker\_results}. Without a reducer, the graph applies updates as
last-write-wins: whichever worker finishes last overwrites every other
worker's output entirely. The naive workaround---giving each worker its
own top-level field, \code{worker\_0\_results} through
\code{worker\_4\_results}---bakes the worker count into the schema and
forces every downstream node to enumerate all five fields by name,
requiring a schema change whenever the count changes. The governing
principle is that a field receiving parallel writes must declare an
accumulation function that matches the semantics required by the
downstream consumer.

\code{Annotated[list[T],~operator.add]} is the LangGraph mechanism for
list-accumulation semantics. The second element of the \code{Annotated}
tuple is the reducer: a callable with signature \code{(left:~T,
right:~T)~->~T} where \code{left} is the current accumulated value and
\code{right} is the incoming partial update from one completed task.
For lists, \code{operator.add} is Python's list concatenation
(\code{left + right}): each completed worker appends its
\code{list[SearchResult]} to the growing accumulated list. By the time
Pregel fan-in fires, the field contains the concatenation of all workers'
lists in task-completion order.

\begin{warningbox}
\code{operator.add} has different semantics for different types.
On lists it concatenates; on numbers it sums. A field declared as
\code{Annotated[float,~operator.add]} with two workers each returning
\code{0.9} produces the accumulated value \code{1.8}. Use a plain
(unannotated) field for scalars written by a single node, or define a
custom reducer such as \code{max} for scalars where the highest
value across workers is the correct merge result.
\end{warningbox}

The choice of \code{operator.add} over a score-max deduplication reducer
is deliberate. Reciprocal Rank Fusion---the aggregation algorithm used in
Section~\ref{sec:map-reduce-topo}---requires knowing each result's rank
position within its originating worker's independent list.
Score-max deduplication in the reducer destroys that information: two
workers returning the same document at different rank positions would be
collapsed into one entry before RRF can compare them. Deduplication must
happen \emph{after} RRF in \code{aggregate\_node}, not before it in the
reducer.

\code{SearchResult} must carry a \code{worker\_id} field from the moment
it is defined. This field is not observability metadata---it is the
structural link that allows \code{aggregate\_node} to reconstruct each
worker's independent ranked list from the flat concatenated
\code{worker\_results} list before passing those lists to RRF.

\begin{lstlisting}[language=Python,
  caption={State schema: list-accumulation reducer and per-result
           worker identifier},
  label={lst:reducers}]
# chapter20/types.py
from pydantic import BaseModel

class SearchResult(BaseModel):
    doc_id:    str    # f"{source}::{sha256(content)[:8]}"
    content:   str
    score:     float  # backend-native score; intra-worker sort key
    source:    str    # "qdrant" | "bm25"
    worker_id: int    # originating search_worker instance

# chapter20/state.py
import operator
from typing import Annotated, TypedDict
from chapter20.types import SearchResult

class WorkerInput(TypedDict):
    sub_query:  str
    worker_id:  int

class ResearchState(TypedDict):
    query:          str
    sub_queries:    list[str]
    worker_results: Annotated[list[SearchResult], operator.add]
    fused_ranking:  list[SearchResult]
    synthesis:      str
\end{lstlisting}

Three decisions in Listing~\ref{lst:reducers} deserve explanation.
\code{SearchResult} is a Pydantic \code{BaseModel} rather than a
\code{TypedDict} because it passes through the LangGraph checkpointer:
\code{BaseModel.model\_dump()} produces a JSON-safe dict without custom
serialisation code, whereas a \code{TypedDict} instance requires explicit
handling when stored in a checkpoint. \code{WorkerInput} lives in
\code{state.py} alongside \code{ResearchState} because it is a state
schema---the shape of a worker node's \code{state} argument---not a
domain value type; keeping both schemas in one file makes the boundary
between parent-graph state and worker-local state immediately visible.
\code{sub\_queries} carries no reducer annotation because it is written
once by \code{decompose\_node} and only read thereafter; declaring a
reducer on a single-writer field adds no value and introduces an
unnecessary merge cost on every checkpoint write.

The reducer contract is now established. The remaining problem is the
fan-out mechanism itself: how to produce a variable number of concurrent
worker invocations whose count is determined at runtime by the query
decomposition result, not at graph-definition time.


%% ------------------------------------------------------------
\section{Dynamic Fan-Out Requires Runtime-Determined Edges}
\label{sec:send-fanout}
%% ------------------------------------------------------------

The design problem is edges whose count is not known at graph-definition
time. \code{add\_conditional\_edges} with a string return value selects
among a fixed set of named destinations; it cannot spawn multiple
concurrent instances of the same node with different inputs. The naive
workaround---five static edges to five statically named worker nodes---
makes worker count a structural constant of the graph and requires each
worker node to know its own index. The governing principle is that the
number and input of parallel tasks must be determined from the state that
exists at the moment fan-out fires, not from the graph's static definition.

\code{Send(node,~arg)} is the LangGraph primitive that closes this gap.
When a conditional edge routing function returns a \code{list[Send]},
LangGraph spawns one independent concurrent task per \code{Send} object.
Each task runs the named node with \code{arg} as its \emph{entire}
state---not as a partial update merged into \code{ResearchState}. The node
receiving the \code{Send} therefore annotates its \code{state} parameter
as \code{WorkerInput}, not \code{ResearchState}: it sees only the two
fields packed into the \code{Send}'s \code{arg} dict. This is the critical
distinction from a normal edge: a normal edge passes the full parent state
to the next node; \code{Send} passes only what the caller explicitly
provides.

\begin{warningbox}
The fan-out routing function must never be registered with \code{add\_node}.
A routing function returns \code{list[Send]}; a node returns a state-update
\code{dict} or \code{Command}. Passing the router to \code{add\_node} causes
LangGraph to attempt to merge the returned \code{list[Send]} into
\code{ResearchState} as a state update, raising a runtime error on the
first invocation. The router appears only as the second argument to
\code{add\_conditional\_edges}.
\end{warningbox}

The decomposition step and the routing step are two units with
two separate contracts. \code{decompose\_node} is an \code{async} node
because it calls an LLM; it writes \code{sub\_queries} to
\code{ResearchState} and returns a state-update \code{dict}.
\code{fan\_out\_router} is a pure synchronous function because it makes
no I/O calls; it reads \code{sub\_queries} from the updated state and
returns \code{list[Send]}. Collapsing them into one function would require
that function to be simultaneously a node (returning a state update) and a
router (returning \code{list[Send]}), which is not a valid LangGraph node
contract.

\begin{lstlisting}[language=Python,
  caption={Decomposition node and pure fan-out routing function},
  label={lst:send-fanout}]
# chapter20/fan_out.py
from langgraph.types import Send
from langchain_core.runnables import RunnableConfig
from langchain_litellm import ChatLiteLLM
from pydantic import BaseModel
from chapter20.state import ResearchState

MAX_WORKERS: int = 5

class SubQueryList(BaseModel):
    queries: list[str]

_llm = ChatLiteLLM(model="anthropic/claude-sonnet-4-8")
_structured_llm = _llm.with_structured_output(SubQueryList)

async def decompose_node(
    state: ResearchState,
    config: RunnableConfig,
) -> dict:
    result: SubQueryList = await _structured_llm.ainvoke(
        f"Decompose into at most {MAX_WORKERS} independent "
        f"sub-queries: {state['query']}"
    )
    return {"sub_queries": result.queries[:MAX_WORKERS]}

def fan_out_router(state: ResearchState) -> list[Send]:
    """Pure routing function — only used in add_conditional_edges."""
    return [
        Send("search_worker", {"sub_query": q, "worker_id": i})
        for i, q in enumerate(state["sub_queries"])
    ]
\end{lstlisting}

Three decisions in Listing~\ref{lst:send-fanout} deserve explanation.
\code{\_structured\_llm} is bound at module scope rather than rebuilt
inside \code{decompose\_node} on each call: \code{with\_structured\_output}
constructs a tool-calling chain whose assembly cost is non-trivial, and
rebuilding it per invocation adds measurable overhead at scale---the same
pattern established for \code{\_structured\_llm} in Chapter~19~\S19.1.
The \code{result.queries[:MAX\_WORKERS]} slice is applied before storing
in state, not inside \code{fan\_out\_router}: the \code{sub\_queries}
field should accurately reflect the work that will be dispatched; capping
in the router would mean the field shows more queries than workers were
actually spawned, misleading any downstream node or audit log that reads
\code{sub\_queries}. \code{fan\_out\_router} passes plain \code{dict}
literals to \code{Send} rather than constructing \code{WorkerInput(...)}
instances because \code{TypedDict} instances are plain \code{dict} objects
at runtime; the \code{WorkerInput} annotation exists only for the
static type-checker inspecting \code{search\_worker}'s signature.

The fan-out contract is established. The problem that fan-out creates
immediately is failure propagation: when two retrieval backends run
concurrently inside a worker, a failure in one must not cancel the other,
and a total worker failure must not prevent the fan-in from ever firing.


%% ------------------------------------------------------------
\section{Partial Failure Must Produce Partial Results, Not No Results}
\label{sec:partial-failure}
%% ------------------------------------------------------------

The design problem is error propagation across two independent concurrency
layers. The worker node is one task in a parallel batch managed by
LangGraph's Pregel scheduler; inside that task, two retrieval backend
calls run concurrently via \code{asyncio.gather}. These layers have
different failure semantics and require different containment mechanisms.
Wrapping the entire worker in a single \code{try/except} that returns an
empty list on any failure silently discards both backends' results when
only one fails---a worse outcome than partial results. The governing
principle is that each layer's failures are contained at that layer:
intra-worker backend failures by \code{asyncio.gather} with
\code{return\_exceptions=True}; transient whole-worker failures by
\code{RetryPolicy} registered at \code{add\_node} time.

\code{asyncio.gather(*coros,~return\_exceptions=True)} changes the gather
contract: instead of raising the first exception and cancelling all
remaining coroutines, it returns a mixed list of results and
\code{Exception} instances in the same positional order as the input
coroutines. The caller then inspects each element with
\code{isinstance(r,~Exception)} and decides what to contribute to the
reducer. A worker whose Qdrant backend times out but whose BM25 backend
succeeds still contributes the BM25 results; without
\code{return\_exceptions=True} the timeout would have cancelled the BM25
call and the worker would contribute nothing.

\code{RetryPolicy} operates at a different level entirely. It is
registered at \code{add\_node} time and applies to the whole node
invocation. When \code{search\_worker} raises an unhandled exception---one
not caught by the \code{asyncio.gather} inspection---LangGraph retries the
whole node with exponential backoff before marking the task as permanently
failed. The two mechanisms are complementary, not alternatives:
\code{return\_exceptions=True} handles cases where partial results are
better than nothing; \code{RetryPolicy} handles cases where a transient
infrastructure failure makes the whole worker temporarily unavailable.

\begin{notebox}
\code{RetryPolicy}'s \code{max\_attempts} field counts the initial attempt
plus all retries. The default \code{RetryPolicy()} uses
\code{max\_attempts=3}: one initial attempt and two retries, not three
retries. The default \code{retry\_on} retries on all exceptions. In
production, restrict it to transient infrastructure exceptions such as
\code{httpx.TimeoutException} and \code{ConnectionError} to avoid
retrying on permanent failures like \code{pydantic.ValidationError} or
authentication errors, which will fail identically on every retry.
\end{notebox}

\begin{lstlisting}[language=Python,
  caption={Worker node with two-level error containment},
  label={lst:gather-worker}]
# chapter20/worker.py
import asyncio
from langchain_core.runnables import RunnableConfig
from chapter20.state import WorkerInput
from chapter20.types import SearchResult
from chapter20.backends import qdrant_search, bm25_search

async def search_worker(
    state: WorkerInput,
    config: RunnableConfig,
) -> dict:
    raw = await asyncio.gather(
        qdrant_search(state["sub_query"],
                      worker_id=state["worker_id"]),
        bm25_search(state["sub_query"],
                    worker_id=state["worker_id"]),
        return_exceptions=True,
    )
    merged: list[SearchResult] = []
    for r in raw:
        if isinstance(r, Exception):
            # Log in production; contribute empty list to reducer.
            continue
        merged.extend(r)
    return {"worker_results": merged}
\end{lstlisting}

Two decisions in Listing~\ref{lst:gather-worker} deserve explanation.
The \code{worker\_id} argument is passed directly to both backend
functions rather than post-assigned on the returned results: the backend
functions construct \code{SearchResult} objects and are the correct place
to assign \code{worker\_id}, keeping it co-located with object construction
rather than requiring a post-hoc list comprehension after the gather call.
The return value \code{\{"worker\_results": merged\}} is a partial update
targeting only the one field this worker owns; the Pregel scheduler merges
it into \code{ResearchState} via the \code{operator.add} reducer, and the
worker has no write access to any other \code{ResearchState} field---a bug
in one worker cannot corrupt fields it never touched.

\begin{lstlisting}[language=Python,
  caption={Backend stubs: typed signatures with \texttt{worker\_id} propagation},
  label={lst:backends}]
# chapter20/backends.py
from chapter20.types import SearchResult

async def qdrant_search(
    query:     str,
    top_k:     int = 20,
    worker_id: int = 0,
) -> list[SearchResult]:
    # Production: see Chapter 11 §11.2 for Qdrant implementation.
    raise NotImplementedError

async def bm25_search(
    query:     str,
    top_k:     int = 20,
    worker_id: int = 0,
) -> list[SearchResult]:
    # Production: see Chapter 11 §11.3 for BM25Okapi implementation.
    raise NotImplementedError
\end{lstlisting}

Two decisions in Listing~\ref{lst:backends} deserve explanation.
Both functions raise \code{NotImplementedError} rather than returning an
empty list: the integration test in Section~\ref{sec:full-parallel} mocks
both backends with fixture data, and a stub that silently returned an
empty list would make tests pass without exercising retrieval, hiding
missing implementations entirely. The \code{worker\_id} default value of
\code{0} allows both functions to be called in backend-level unit tests
without constructing a full \code{WorkerInput} context.

The worker's two-level error containment is correct. What remains is the
structural question: how does Pregel fan-in actually fire, what does
\code{aggregate\_node} receive when it runs, and how does it transform a
flat accumulated list into a ranked output ready for synthesis.


%% ------------------------------------------------------------
\section{Fan-In Topology: One Aggregation Node Receives the Full Merge}
\label{sec:map-reduce-topo}
%% ------------------------------------------------------------

The design problem is the contract between parallel workers and their
aggregation step. The edge declaration
\code{add\_edge("search\_worker",~"aggregate")} reads like a per-worker
trigger, which invites designs where the aggregation node checks how many
workers have reported and accumulates state internally across repeated
calls. That design is incorrect and unnecessary: it requires knowing the
expected worker count at aggregation time, which changes as
\code{sub\_queries} changes between runs. The governing principle is that
LangGraph's Pregel scheduler is the accumulator: \code{aggregate\_node}
runs exactly once, after all tasks in the current \code{Send} batch have
completed, and receives \code{ResearchState} with \code{worker\_results}
already fully merged by the \code{operator.add} reducer.

The fan-in mechanism is implicit in one edge declaration. As each parallel
task completes, its \code{\{"worker\_results":~...\}} partial update is
applied through the \code{operator.add} reducer to the accumulated state.
When the last task in the batch finishes---or is permanently failed after
\code{RetryPolicy} exhaustion---the graph advances to \code{"aggregate"}.
The aggregation node never checks a worker count or tests completeness:
the scheduler guarantees both conditions before the node runs.

RRF requires per-worker ranked lists, but \code{operator.add} delivers a
flat concatenated list. The bridge is the \code{worker\_id} field on every
\code{SearchResult}: grouping by \code{worker\_id} reconstructs each
worker's independent list that the reducer intentionally preserved by
never deduplicating. Sorting each group by \code{score} descending
restores intra-worker rank order that task-completion interleaving may
have scrambled. Chapter~11~\S11.4 introduced \code{reciprocal\_rank\_fusion}
as a function that accepts a list of ranked lists and returns a unified
re-ranked list; this chapter imports that implementation without
modification.

\begin{warningbox}
If any \code{search\_worker} task raises an unhandled exception after
\code{RetryPolicy} exhaustion, LangGraph marks that task as permanently
failed. Whether Pregel fan-in then fires depends on the graph's configured
error policy. To guarantee \code{aggregate\_node} always runs even when
some workers fail permanently, workers must catch all exceptions and
return \code{\{"worker\_results":[]\}} rather than raising---exactly what
the \code{isinstance(r,~Exception)} guard in
Listing~\ref{lst:gather-worker} achieves by using \code{continue} instead
of \code{raise}.
\end{warningbox}

\begin{lstlisting}[language=Python,
  caption={Aggregation node: per-worker grouping, RRF, post-fusion
           deduplication},
  label={lst:aggregate}]
# chapter20/aggregate.py
from langchain_core.runnables import RunnableConfig
from chapter20.state import ResearchState
from chapter20.types import SearchResult
from chapter11.fusion import reciprocal_rank_fusion

async def aggregate_node(
    state: ResearchState,
    config: RunnableConfig,
) -> dict:
    by_worker: dict[int, list[SearchResult]] = {}
    for r in state["worker_results"]:
        by_worker.setdefault(r.worker_id, []).append(r)
    ranked_lists = [
        sorted(recs, key=lambda x: x.score, reverse=True)
        for recs in by_worker.values()
    ]
    fused: list[SearchResult] = reciprocal_rank_fusion(
        ranked_lists, k=60,
    )
    deduped: list[SearchResult] = []
    seen: set[str] = set()
    for r in fused:
        if r.doc_id not in seen:
            seen.add(r.doc_id)
            deduped.append(r)
    return {"fused_ranking": deduped}
\end{lstlisting}

Three decisions in Listing~\ref{lst:aggregate} deserve explanation.
Deduplication by \code{doc\_id} happens after RRF, not before it in the
reducer: RRF's cross-worker rank comparison requires seeing the same
document at different rank positions across different workers' lists;
removing duplicates earlier would collapse that signal.
The deduplication uses an explicit \code{for} loop with a \code{seen} set
rather than a one-liner list comprehension that exploits
\code{set.add} returning \code{None}: the loop makes the first-occurrence
selection intent unambiguous and passes \code{mypy} without a
\code{\#~type:~ignore} comment. The \code{by\_worker} grouping and
intra-group sort are needed even though the reducer preserved all records
with their \code{worker\_id} intact, because \code{operator.add}
interleaves results from different workers in task-completion order---a
fast-finishing worker's results may appear at positions that no longer
reflect their score rank within that worker's contribution.

\begin{lstlisting}[language=Python,
  caption={Synthesis node: top-ten fused results as numbered snippets},
  label={lst:synthesise}]
# chapter20/aggregate.py  (continued)
from langchain_litellm import ChatLiteLLM

_synth_llm = ChatLiteLLM(model="anthropic/claude-sonnet-4-8")

async def synthesise_node(
    state: ResearchState,
    config: RunnableConfig,
) -> dict:
    snippets = "\n\n".join(
        f"[{i + 1}] {r.content}"
        for i, r in enumerate(state["fused_ranking"][:10])
    )
    prompt = (
        f"Query: {state['query']}\n\n"
        f"Sources:\n{snippets}\n\n"
        "Write a synthesised answer. "
        "Cite sources inline as [1], [2], etc."
    )
    response = await _synth_llm.ainvoke(prompt)
    return {"synthesis": response.content}
\end{lstlisting}

Two decisions in Listing~\ref{lst:synthesise} deserve explanation.
\code{synthesise\_node} is a separate node from \code{aggregate\_node}
rather than a second phase inside the same function: aggregation is purely
deterministic (sorting and RRF always produce the same output for the same
input), making it a candidate for \code{CachePolicy} in
Section~\ref{sec:full-parallel}; synthesis calls an LLM and is
non-deterministic, making caching incorrect for it. Mixing the two in one
node would prevent applying \code{CachePolicy} to either step without
affecting the other. The \code{\_synth\_llm} binding is declared
immediately before \code{synthesise\_node} in the file rather than at the
top of \code{aggregate.py}: a module-level LLM client is a meaningful
side effect at import time, and positioning it next to its sole consumer
makes the dependency visible without scrolling past unrelated code.

All five nodes are individually defined and consistent in their type
contracts. The remaining problem is assembly: wiring them into a graph
that correctly sources the conditional edge, applies \code{RetryPolicy}
and \code{CachePolicy} at node registration, and is verified by an
integration test that mocks backends rather than calling unreachable stubs.


%% ------------------------------------------------------------
\section{Full System: Five Concurrent Searchers Under One Coordinator}
\label{sec:full-parallel}
%% ------------------------------------------------------------

The design problem is silent topology divergence between the intended
graph structure and the compiled one. Three assembly mistakes are possible
in this topology, none of which produce a compile-time error. First:
registering \code{fan\_out\_router} with \code{add\_node}, which causes
LangGraph to treat \code{list[Send]} as a state update. Second: sourcing
\code{add\_conditional\_edges} from a non-existent node name such as
\code{"fan\_out"} instead of \code{"decompose"}. Third: omitting the
\code{["search\_worker"]} path-map argument from
\code{add\_conditional\_edges}, which removes compile-time validation that
all \code{Send} destinations exist as registered nodes. The governing
principle is that the graph assembly function must make topology explicit
in a form that can be read and reviewed without running the graph.

\code{RetryPolicy} and \code{CachePolicy} are passed as keyword arguments
to \code{add\_node}, not to \code{compile} or to the node function itself.
\code{RetryPolicy()} with default arguments is applied to
\code{"search\_worker"} because workers call external backends subject to
transient failures. \code{CachePolicy()} with default arguments is applied
to \code{"aggregate"} because aggregation is deterministic: given the same
\code{worker\_results} input, RRF and deduplication always produce the
same \code{fused\_ranking}. LangGraph uses a pickle-hash of the node's
input as the default cache key, meaning any change in
\code{worker\_results}---including order differences from non-deterministic
task completion---produces a cache miss.

\begin{notebox}
\code{CachePolicy(ttl=None)} means entries never expire within the
process. With \code{MemorySaver}, both the checkpoint store and the cache
are in-process memory, so a restart clears both regardless of
\code{ttl}; the setting has no practical effect in development.
\code{ttl} becomes meaningful only with a persistent checkpointer such as
\code{AsyncSqliteSaver}: a cached aggregation from a previous run with
identical \code{worker\_results} would survive a process restart and be
served from cache, saving the RRF computation entirely. Chapter~22~\S22.1
introduces \code{AsyncSqliteSaver} and the point at which \code{ttl}
should be set to a bounded value for production deployments.
\end{notebox}

\begin{lstlisting}[language=Python,
  caption={Graph assembly: correct fan-out topology, \texttt{RetryPolicy},
           and \texttt{CachePolicy}},
  label={lst:full-graph}]
# chapter20/graph.py
from langgraph.graph import StateGraph, START, END
from langgraph.graph.state import CompiledStateGraph
from langgraph.types import RetryPolicy, CachePolicy
from langgraph.checkpoint.memory import MemorySaver
from langchain_core.runnables import RunnableConfig
from chapter20.state import ResearchState
from chapter20.fan_out import decompose_node, fan_out_router
from chapter20.worker import search_worker
from chapter20.aggregate import aggregate_node, synthesise_node

def build_research_graph() -> CompiledStateGraph:
    b = StateGraph(ResearchState)
    b.add_node("decompose",     decompose_node)
    b.add_node("search_worker", search_worker,
               retry=RetryPolicy())
    b.add_node("aggregate",     aggregate_node,
               cache=CachePolicy())
    b.add_node("synthesise",    synthesise_node)
    b.add_edge(START, "decompose")
    b.add_conditional_edges(
        "decompose", fan_out_router, ["search_worker"],
    )
    b.add_edge("search_worker", "aggregate")
    b.add_edge("aggregate",     "synthesise")
    b.add_edge("synthesise",    END)
    return b.compile(checkpointer=MemorySaver())

async def run_research(query: str, thread_id: str) -> ResearchState:
    graph = build_research_graph()
    config: RunnableConfig = {
        "configurable": {"thread_id": thread_id},
    }
    return await graph.ainvoke(
        ResearchState(
            query=query, sub_queries=[],
            worker_results=[], fused_ranking=[], synthesis="",
        ),
        config=config,
    )
\end{lstlisting}

Three decisions in Listing~\ref{lst:full-graph} deserve explanation.
\code{add\_conditional\_edges("decompose",~fan\_out\_router,~["search\_worker"])}
sources the conditional edge from \code{"decompose"}---the last node that
ran before routing is needed, and the node that wrote \code{sub\_queries}
to state---not from a hypothetical \code{"fan\_out"} node that is not
registered. The \code{["search\_worker"]} list is the path map: it tells
LangGraph which node names are valid destinations for this conditional edge,
enabling compile-time validation that every name produced by a
\code{Send} in \code{fan\_out\_router} is a registered node; without it, a
typo in \code{fan\_out\_router}'s \code{Send} node name produces a
\code{NodeNotFoundError} only at runtime on the first query.
\code{RetryPolicy()} is instantiated with no arguments to make explicit that
all defaults are intentional; passing \code{RetryPolicy(max\_attempts=3)}
would suggest a deliberate override from a different default, when
\code{max\_attempts=3} is in fact the default value.

\begin{lstlisting}[language=Python,
  caption={Integration test with backend mocks and structural assertions},
  label={lst:full-test}]
# chapter20/test_graph.py
import pytest
from unittest.mock import AsyncMock, patch
from chapter20.graph import run_research
from chapter20.fan_out import MAX_WORKERS
from chapter20.types import SearchResult

def _make_results(worker_id: int, n: int = 3) -> list[SearchResult]:
    return [
        SearchResult(
            doc_id=f"w{worker_id}::doc{i}",
            content=f"Content {i}",
            score=1.0 / (i + 1),
            source="mock",
            worker_id=worker_id,
        )
        for i in range(n)
    ]

@pytest.mark.asyncio
async def test_parallel_research_produces_synthesis() -> None:
    mock_r = _make_results(worker_id=0)
    with patch("chapter20.backends.qdrant_search",
               new=AsyncMock(return_value=mock_r)), \
         patch("chapter20.backends.bm25_search",
               new=AsyncMock(return_value=mock_r)):
        result = await run_research(
            query="Quantum error correction methods.",
            thread_id="test-ch20-001",
        )
    assert result["sub_queries"], \
        "decompose_node must produce at least one sub-query"
    assert result["fused_ranking"], \
        "aggregate_node must produce at least one fused result"
    assert result["synthesis"] != "", \
        "synthesise_node must write a non-empty synthesis"

@pytest.mark.asyncio
async def test_worker_count_bounded_by_max_workers() -> None:
    mock_r = _make_results(worker_id=0)
    with patch("chapter20.backends.qdrant_search",
               new=AsyncMock(return_value=mock_r)), \
         patch("chapter20.backends.bm25_search",
               new=AsyncMock(return_value=mock_r)):
        result = await run_research(
            query="A highly compound query " * 10,
            thread_id="test-ch20-002",
        )
    assert len(result["sub_queries"]) <= MAX_WORKERS

@pytest.mark.asyncio
async def test_partial_backend_failure_still_synthesises() -> None:
    pytest.skip(
        "Requires fault-injection fixture: Chapter 22 §22.4"
    )
\end{lstlisting}

Three decisions in Listing~\ref{lst:full-test} deserve explanation.
\code{\_make\_results} is a module-level fixture factory rather than inline
dict construction: it ensures every test that needs \code{SearchResult}
objects uses the same valid \code{doc\_id} format
(\code{f"w\{worker\_id\}::doc\{i\}"}), making the format contract testable
rather than repeated and divergent across test functions. The \code{patch}
calls target the names in \code{chapter20.backends}---the module where
they are defined and imported from---rather than at their reference site in
\code{chapter20.worker}: Python's \code{unittest.mock.patch} replaces the
name in the namespace it is given; patching at the wrong module leaves the
original unpatched and the test silently calls the \code{NotImplementedError}
stubs. \code{test\_partial\_backend\_failure\_still\_synthesises} calls
\code{pytest.skip} with an explicit message naming the missing capability
rather than an empty body: the test's name precisely states the invariant it
will verify once the fault-injection fixture exists, making it a visible
specification item in the suite rather than dead code.

\begin{notebox}
The parallel research graph as built here dispatches all sub-queries to the
same \code{search\_worker} node. When sub-queries differ substantially in
domain---one requiring dense vector retrieval, another requiring structured
knowledge-graph traversal---a single worker type produces uniformly mediocre
results. A later chapter shows how to combine \code{Send}-based fan-out
with the Chapter~19 supervisor pattern to route different sub-queries to
different specialised worker types while preserving the fan-in topology and
reducer contracts built here.
\end{notebox}

Chapter~21 introduces a new dimension to the execution model built here:
where this chapter parallelises work that has no data dependency and no
need for human review, the next chapter confronts the case where a graph
must pause mid-execution---before fan-out commits five concurrent retrieval
calls---and surface the decomposed sub-queries to a human reviewer who may
edit, reject, or approve them before execution resumes.
